\chapter{Qué habría que hacer}

\label{cap:plan}
\begin{cautela}
Este capítulo es, junto con el apéndice \ref{cap:propuestas}, el único lugar del libro donde se propone algo, y conviene decir con
precisión qué es y qué no es.

Va como capítulo y no como apéndice porque es prosa y no aparato, y porque sigue
directamente al anterior: aquel describe qué es este territorio; éste, qué podría hacerse.
Pero conviene que el lector lo separe con nitidez de los veintisiete capítulos que lo preceden.

\textbf{No es un plan técnico.} Este trabajo sostiene, capítulo tras capítulo, que faltan los
datos con los que se planifica: no se afora el río desde 1986 --- las cinco estaciones de la cuenca están cerradas y sus series no alcanzan el mínimo que la propia norma exige ---, no se conoce el antecedente dominial de las matrículas centrales, no se publicaron las cotas de inundación ni el padrón de servicios, y el ancho de la franja no urbanizable del río no está expresado en metros en ninguna parte del Código que la establece. Proponer obras concretas sobre esa base sería incurrir en lo mismo que el libro
señala.

\textbf{Es un plan institucional.} Enumera qué hay que medir antes de decidir, qué hay que
decidir, quién tiene competencia para hacerlo y en qué orden. Nada de lo que sigue requiere
información que este libro no tenga: requiere decisiones.

Y tiene un fundamento que conviene enunciar de entrada, porque es lo que lo separa de una
lista de deseos. \textbf{Casi todo lo que se propone aquí ya se hizo alguna vez en este mismo
departamento}, y el cuerpo del libro lo documenta. No hay que inventar instrumentos: hay que
volver a usar los que existen.
\end{cautela}

\section{Primero: lo que no se puede planificar sin medir}

Cuatro mediciones faltan, y sin ellas cualquier plan es una apuesta. Las cuatro son
producibles en meses, no en años, y ninguna requiere legislación nueva.

\textbf{Aforar los ríos.} En 1924 el Departamento de Obras Públicas dejó constancia escrita
de que no podía informar sobre el caudal del río Wierna «en razón de no haberse efectuado
aforos de los ríos de la Provincia». En 2015, el procedimiento de línea de ribera seguía
poniendo los aforos en primer lugar de prelación --- con series de cincuenta años --- y
resolviéndose en la práctica por método geológico-geomorfológico. \textbf{Cien años después, la cuenca no tiene una sola estación de aforo en funcionamiento: la última cerró en 1986.}

Sin aforos no se sabe cuánta agua hay, cuánta se está concediendo ni cuánta queda. El
capítulo \ref{cap:loteo} documenta tres gestiones de abastecimiento poblacional con tratamientos distintos
y sin fundamento explícito; el capítulo \ref{cap:politica} documenta que el mismo desconocimiento sirvió, en
1936, para negarle una indemnización a quien había ganado su juicio. \textbf{La ausencia de
aforos no es neutral: favorece a quien decide, porque vuelve indiscutible cualquier decisión.}

\textbf{Publicar las cotas y las líneas de ribera con sus coordenadas.} Doce de quince determinaciones del período relevado publicaron sus coordenadas; de las tres que no lo hicieron, las dos referidas a un inmueble recaen sobre suelo en proceso de loteo, condición que comparten otras dos que sí las publicaron. Publicarlas todas es un acto administrativo de una página, y hacerlas legibles ---las del pueblo están sólo en una imagen escaneada--- es otro.

\textbf{Levantar el estado dominial de las fincas del corredor.} El departamento está
catastrado desde 1929. El registro existe, funciona y no es consultable sin arancel ni
presencia física. Un informe de dominio de las matrículas 4.366 y 6.558 --- las dos que
atraviesan este libro --- cuesta lo que cuesta un trámite y contestaría la pregunta central
del capítulo \ref{cap:loteo}.

\textbf{Publicar el padrón de servicios.} Qué manzanas tienen red de agua, cuáles cloaca,
cuáles gas, y cuáles no. Sin ese mapa no se puede saber si un loteo nuevo está donde hay
infraestructura o donde no la hay, que es exactamente lo que la Ordenanza 437 de Vaqueros se
propuso evitar en 2015.

\section{El río}

\subsection*{Lo que el propio Estado ya diagnosticó}

En diciembre de 1935, la Dirección General de Obras Públicas elevó un informe que sigue
siendo la mejor descripción disponible del problema. Identificó que el río se bifurca en tres
brazos cuatro kilómetros aguas arriba del pueblo, que el de la margen derecha se lleva la
mayor parte del caudal y que toma una curva brusca al pasar junto a las casas. Y propuso
intervenir \textbf{en la bifurcación} y no junto a las viviendas.

Ese informe agregó algo más: dijo que las obras que autorizaba eran provisorias y servirían
para «adquirir experiencia sobre el comportamiento del Río», de modo que después «se puedan
proyectar obras más duraderas».

\textbf{Ese estudio nunca se hizo.} Noventa y un años después, lo que se sigue construyendo
son espigones de emergencia, bateas y muros de contención.

\subsection*{Qué habría que hacer}

\textbf{Encargar el estudio que falta desde 1935.} Un estudio hidrológico e hidráulico de la
cuenca completa, con una serie de aforos que hoy no se produce y modelización de crecidas. Es la
condición de todo lo demás, y es lo que la sentencia de 2021 ordena cuando manda privilegiar
infraestructura verde por sobre la gris: no se puede decidir qué obra corresponde sin saber
cómo se comporta el río.

\textbf{Crear una autoridad de cuenca.} El río La Caldera nace en un municipio, atraviesa
otro, cambia de nombre y desemboca en un tercer departamento. El capítulo \ref{cap:redes} documenta que en
1929 la Provincia resolvió ese problema \textbf{repartiendo el río}: le dio a la
Municipalidad de Campo Santo potestades sobre el tramo que le importaba, incluida la de
prohibir nuevas tomas. Noventa y siete años después, ninguna jurisdicción administra el
conjunto del cauce, la extracción de áridos se concesiona en ambas márgenes sin autoridad
común, y la única figura metropolitana que la provincia construyó para el corredor es la del
transporte de pasajeros.

La Ley 7322 demuestra que la Provincia sabe crear autoridades metropolitanas cuando decide
hacerlo. Que exista una para los colectivos y ninguna para la cuenca no es un problema
constitucional: es una decisión sobre qué se administra en escala metropolitana y qué no.

\textbf{Y evaluar el efecto acumulado de un siglo de intervenciones.} Este libro documenta
obras sobre el mismo cauce en 1910, 1911, 1935, 1936, 1938, 1940, 1945, 1951, 1980, 2001, 2010, 2014, 2015, 2016--2017, 2023 y 2026, por
distintos organismos, sin que ningún acto evalúe el conjunto. En 1940 el propio Estado
consignó que una obra proyectada hubo de suspenderse \textbf{porque el río desvió su curso}.
El artículo 171 del Código de Aguas exige evaluación previa para los permisos de extracción
de áridos; para las obras del propio Estado sobre el cauce no se encontró norma equivalente.
Establecerla es una decisión de la autoridad de aplicación.

\section{El suelo}

\subsection*{Los loteos que ya existen}

Nada de lo que sigue propone desconocer derechos adquiridos, y conviene decirlo: quien compró
un lote de buena fe no es responsable de las omisiones del procedimiento que lo aprobó.

\textbf{Determinar qué está dentro de la línea de ribera y qué no.} Es información que en
buena parte ya existe: veinticinco actos entre 2013 y 2015 sobre más de veinticinco inmuebles.
Publicarla en un plano único y accesible convierte veinticinco actos dispersos en un
instrumento de gestión.

\textbf{Completar los expedientes incompletos.} El plano de El Durazno se aprobó sólo por las
etapas 1, 2 y 4; las etapas 3 y 5 no constan. Ese vacío no protege a nadie: no protege al
municipio, que no sabe qué habilitó, ni al comprador, que no sabe si su lote está aprobado.
El vecino de La Calderilla que en enero de 2026 preguntó al municipio si los loteos estaban
aprobados y no recibió respuesta está describiendo un problema con solución administrativa
simple.

\textbf{Exigir el cumplimiento de las obras comprometidas.} El Código de Ordenamiento
Territorial prevé Garantía de Ejecución de Obras y actas de Control de Avance. Aplicarlos no
requiere norma nueva. El anexo del Decreto 1682/19 prevé, además, caución por el monto de las obras y su ejecución si el desarrollador no las termina.

\subsection*{Los loteos futuros}

\textbf{Aplicar las condiciones que el propio Estado ya supo exigir.} El capítulo \ref{cap:loteo} documenta
que a Las Vertientes se le exigió resolución de línea de ribera, y que la línea se determinó
en sesenta y siete días; que el plano del Barrio La Misión descuenta la superficie del retiro
de ribera y declara reserva no urbanizable; y que a Chacras del Gallinato se le exigieron seis
condiciones, incluido un decreto del gobernador otorgando la concesión de agua.

Ninguna de esas exigencias es una invención: las tres son actos de la cartera provincial que
aprueba los loteos ---entonces, la de Economía---, sobre este mismo departamento, entre 2014 y 2015. \textbf{Que El Durazno recibiera dos
condiciones en 2021 no se explica por imposibilidad técnica.}

\textbf{Aplicar el artículo 172 del Código de Aguas.} Prohíbe instalar asentamientos poblacionales en los sectores que la Autoridad determine como inundables, y exige autorización previa ---con obras de prevención--- para toda construcción en zona ribereña. Está vigente y es de aplicación
obligatoria.

\textbf{Y usar la audiencia pública ambiental.} El artículo 49 de la Ley 7070 se aplicó en
2014 a una urbanización del departamento: la audiencia se celebró en el Centro Deportivo
Municipal de Vaqueros y el expediente completo quedó disponible en el propio municipio. El
Estado sabe acercar los expedientes cuando decide hacerlo.

\section{Agua, cloaca y gas}

\subsection*{La secuencia está invertida}

El problema de fondo no es la falta de proyectos. Este libro documenta que en 2014 la
Provincia aprobó el proyecto ejecutivo del acueducto de Vaqueros ---confeccionado por la prestadora--- y autorizó el concurso de la obra, y que en 2015 licitó redes
colectoras, colectora máxima y planta depuradora cloacal para esa localidad. Los proyectos
existieron.

Lo que documenta también es que \textbf{en los mismos años el Estado homologaba
abastecimientos privados}: a un club de campo se le exigió prever el servicio de agua «en
forma privada»; otro gestionó ciento quince metros cúbicos diarios de un río con
carácter permanente; y para cuatrocientos lotes de un tercero se tramitó un dren horizontal.

\textbf{La secuencia correcta es la que la propia Ordenanza 437 de Vaqueros enunció en
enero de 2015}: «lograr el máximo aprovechamiento de la infraestructura existente,
evitando la habilitación de nuevas áreas sin disponibilidades de extensión de las mismas». No
habilitar suelo donde no se pueda llevar servicios. El reglamento que fija cómo se aplica ese
criterio es el anexo que el Boletín nunca publicó y que el municipio terminó ofreciendo en su digesto en línea: exige agua corriente, línea de ribera, cesiones y garantía antes de vender (capítulo \ref{cap:vaqueros}).

\subsection*{Qué habría que hacer}

\textbf{Publicar el estado de ejecución de lo ya licitado.} El expediente 7807/15 de la obra
cloacal de Vaqueros: si se adjudicó, si se ejecutó, qué trazado tuvo y qué proporción del
ejido quedó servida. Es una consulta, no una obra.

\textbf{Resolver el problema de las servidumbres, que es el nudo real.} La circular
aclaratoria de esa licitación revela que la red debía atravesar propiedad privada en todos los
tramos proyectados sobre terrenos privados, y que el estudio de títulos y los planos de
servidumbre quedaban a cargo del contratista. \textbf{Cuando la infraestructura pública
depende de negociar paso por fincas privadas, el trazado sigue el mapa de la propiedad y no
el de la población.} Constituir esas servidumbres por vía administrativa, con la Dirección
General de Inmuebles interviniendo de oficio, es una decisión que cambia el problema de
naturaleza.

\textbf{Condicionar la aprobación de nuevo suelo a la factibilidad de servicios.} No es una
propuesta original: es lo que dice el considerando de una ordenanza vigente desde 2015.

\textbf{Pedir el cálculo que nadie pidió para el gas.} La red llega a Vaqueros y se detiene allí (capítulo \ref{cap:redes}). El régimen de expansión de ENARGAS permite que terceros interesados ---el municipio, un grupo de vecinos--- soliciten a la distribuidora el cálculo del Valor de Negocio de una extensión hacia La Caldera. Ese cálculo dice si la obra le corresponde a la distribuidora o cuánto deben aportar los usuarios. Pedirlo no cuesta nada y convierte una carencia en un número.

\textbf{Y decidir, explícitamente, si el abastecimiento privado es una solución transitoria o un régimen.} Hoy conviven las dos cosas sin que ningún acto las distinga. Un club de campo con
concesión permanente sobre un río y un vecino de la zona alta con camión cisterna no son dos
etapas de un mismo proceso: son dos regímenes distintos aplicados al mismo territorio.

\section{La capacidad institucional}

Nada de lo anterior se sostiene sin esto, y es lo más difícil.

Este libro documenta que el municipio de La Caldera fue clasificado en tercera categoría
el 9 de diciembre de 1931 y que la estadística provincial de 2022 lo sigue clasificando así, como a Vaqueros, aunque la ley que definía las categorías se derogó en 2018; que en 1908 vendía en licitación el derecho a cobrar sus
impuestos; que en 1922 el Interventor Nacional designaba a los cinco miembros de su comisión y le
aprobaba el presupuesto; que en 1933 gastaba en sueldos administrativos más de lo que la ley
permitía; que en 1940 debía a Sanidad cerca de un tercio de su renta anual; y que en julio de
2020 declaró por escrito no poder pagar los sueldos de su personal.

\textbf{Dejar de describir a los municipios con una categoría que ya no existe} y medirlos por lo que la Ley 8126 les exige a todos por igual ---Concejo, intendente, plan de obras, presupuesto con límite de gasto en personal, cuenta general ante la Auditoría, publicación de sus normas--- es el primer paso para dimensionar lo que les falta. La clasificación de 1931 era
correcta cuando se hizo; lo que nunca ocurrió es que volviera a hacerse, y la ley que obligaba a rehacerla ya no rige.

\textbf{Revisar la base imponible.} El capítulo \ref{cap:hacienda} documenta lotes urbanos con valuación
fiscal de tres cifras: un municipio cuya superficie urbana se multiplica sobre esas
valuaciones aumenta su obligación de servicio sin aumentar su recaudación. Un revalúo no es
un aumento de impuestos: es la condición para que la superficie loteada financie los
servicios que exige.

\textbf{Y devolverle voz sobre sus recursos.} En 1914 el municipio de La Caldera dictó una
ordenanza que reglamentaba la distribución de la totalidad del caudal de uno de sus ríos,
reconociendo los usos consuetudinarios existentes. En 1939 esa ordenanza prevaleció sobre una
solicitud de concesión provincial, y al año siguiente el Poder Ejecutivo ratificó que no
podía conceder agua de ese río. En 1924, 1938 y 1939 el dictamen o el certificado del
municipio integraron expedientes de agua y quedaron publicados.

Entre 2013 y 2015, en cambio, la Provincia dictó veinticinco actos de línea de ribera sobre la
cuenca de este departamento y en ninguno consta intervención municipal. La excepción, que la serie cuenta aparte, prueba que la práctica no se había perdido del todo: la Resolución 023/14 se dictó a pedido del intendente, sobre la ribera del propio pueblo.

\textbf{Que el municipio sea oído en los trámites hídricos y de fraccionamiento de su propio
territorio no es una innovación: es restituir una práctica que existió y que se perdió.}

\section{Qué de todo esto ya está ordenado}

Conviene cerrar señalando que una parte sustancial de lo anterior no requiere decisión nueva,
porque ya fue ordenada judicialmente.

La sentencia de 2021 declaró la obligación estatal de proteger y recomponer la cuenca, con
criterio ecocéntrico, y ordenó un plan de manejo que privilegiara infraestructura verde por
sobre la gris. El plan se aprobó en febrero de 2026, cuatro años y dos meses después, bajo apercibimiento de remisión a la justicia penal. Este libro lo leyó: cifra sus medidas, proyecta una alerta temprana y asigna menos del uno por ciento de su inversión a infraestructura verde (capítulo \ref{cap:amparo}).

\textbf{El primer paso de cualquier plan es, entonces, partir del que ya existe} y ponerlo a disposición pública más allá de un domicilio particular. El apéndice \ref{cap:pedidos} reclama además la sentencia completa.

\section{Qué de todo esto no está ordenado}

Y conviene decir lo contrario con la misma claridad, porque es lo que separa a este capítulo
de una lista de deseos.

\textbf{Una sentencia puede ordenar un plan. No puede ordenar un criterio.} Todo lo que sigue
es legal, está vigente y no lo alcanza ninguna orden judicial ni ninguna norma que obligue a
cambiarlo. Si alguien quisiera modificarlo, tendría que decidirlo.

\textbf{Nadie está obligado a publicar el atlas del Código en un formato que sirva.} Los ocho planos que definen dónde se puede construir y dónde no integran la ordenanza por su artículo 12, y se publicaron sólo como imágenes a escala municipal, sin las capas georreferenciadas con que se confeccionaron. No
hay plazo, no hay sanción y no hay un solo acto que lo requiera.

\textbf{Nadie está obligado a decir cuántos metros mide la franja.} El capítulo
\ref{cap:cot} lo buscó en los tres lugares donde podía estar y no está. El Código remite a un
plano, el plano está a escala de cuatro kilómetros, y ninguna norma exige traducir el trazo a
una cifra.

\textbf{Nadie está obligado a explicar por qué unas coordenadas se publican y otras no.} El
capítulo \ref{cap:ribera} encontró dos determinaciones sobre la misma matrícula, del mismo
organismo, separadas por cinco meses: una publica ciento trece pares de coordenadas y la
otra las remite a la sede. Ninguna norma manda una cosa ni la otra, y por lo tanto ninguna
obliga a fundar la diferencia.

\textbf{Nadie está obligado a resolver las contradicciones de su propio código.} El artículo
12 dice que cuando el texto y el plano se contradicen prevalece lo que determine el Ejecutivo
Municipal, previo dictamen. No dice que el dictamen deba publicarse.

\textbf{Nadie está obligado a que el Boletín traiga lo que anuncia.} La edición que publica el
decreto que rige los veinticinco actos de línea de ribera de este departamento trae una línea que remite
a una separata, y la separata no está digitalizada. Es un instrumento previsto y legítimo.

\textbf{Nadie está obligado a conciliar dos registros del mismo Estado.} Veinte concesiones de
áridos figuran como vigentes; seis de ellas, con la evaluación ambiental expresamente no
aprobada. Son dos registros de dos organismos, y ninguna norma manda cruzarlos.

\textbf{Nadie obligó a un municipio a contestar.} En la ejecución de la sentencia, uno de los
dos municipios de la cuenca no presentó el informe que se le requirió. El expediente lo
consigna. Este libro no encontró consecuencia alguna.

\textbf{Y nadie está obligado a que los dos circuitos se parezcan.} El Estado adquiere por
decreto un inmueble privado que ocupa hace cuarenta y siete años; un vecino que ocupa un inmueble privado el mismo tiempo litiga, y el que ocupa tierra fiscal espera, a veces treinta años, un decreto. Las dos cosas están en la ley, y son leyes distintas.

\begin{cautela}
Ninguno de estos ocho puntos describe una ilegalidad. Todos describen \textbf{un margen}: un
lugar donde el ordenamiento deja que se decida y no exige decir por qué. Sumados, son la
segunda tesis de este libro en su forma más concreta.

Por eso la lista importa para un plan. \textbf{Lo que una sentencia puede ordenar, ya fue
ordenado.} Lo que queda es aquello que ninguna sentencia alcanza, porque no se trata de
incumplimientos sino de criterios, y los criterios no se ejecutan: se adoptan.

\textbf{Y se adoptan sin el costo de una obra.} Publicar las capas del atlas, informar el ancho
en metros, fundar por escrito por qué una coordenada se publica y otra no, cruzar dos
registros: nada de eso requiere presupuesto. Requiere que alguien decida que de aquí en más se
hace así, y lo escriba.
\end{cautela}

\section{Un orden de prioridades}

Si hubiera que ordenar todo lo anterior por lo que cuesta y por lo que habilita, el orden
sería éste.

\begin{enumerate}[leftmargin=*]
\item \textbf{Publicar lo que ya existe}: el Plan de Manejo aprobado, la sentencia, las coordenadas de las determinaciones que no las publicaron, en formato legible, el estado del expediente cloacal de Vaqueros y el padrón de servicios. Costo:
administrativo. Habilita: todo lo demás.
\item \textbf{Medir}: aforos, estudio hidrológico de cuenca, informe de dominio de las
matrículas centrales, revalúo inmobiliario. Costo: moderado. Plazo: meses.
\item \textbf{Aplicar lo vigente}: el artículo 172 del Código de Aguas, los certificados del
Código de Ordenamiento Territorial, la audiencia pública del artículo 49 de la Ley 7070, el
criterio de infraestructura previa de la Ordenanza 437. Costo: nulo. Requiere sólo decisión.
\item \textbf{Decidir}: la autoridad de cuenca, la recategorización municipal, la
constitución administrativa de servidumbres, el régimen del abastecimiento privado. Costo:
político. Requiere norma provincial.
\item \textbf{Construir}: lo que el estudio hidrológico indique. Es lo último, no lo primero,
y es exactamente lo que este departamento viene haciendo primero desde 1910.
\end{enumerate}

\begin{cautela}
Una última advertencia, que corresponde al carácter de este capítulo.

Este libro no sabe cuánto cuesta nada de esto, ni de dónde saldrían los fondos, ni si existe
la voluntad política de hacerlo. Tampoco sabe si los vecinos del departamento --- que no
fueron entrevistados para este trabajo --- estarían de acuerdo con este orden de prioridades.

Lo que sí puede afirmar, y es todo lo que este capítulo pretende, es que \textbf{ninguno de
los instrumentos necesarios falta}. La determinación de línea de ribera existe y funciona. La
audiencia pública ambiental existe y se aplicó acá. La exigencia de agua previa a la
asignación de matrículas existe y se impuso acá. La autoridad metropolitana existe, para el
transporte. La ordenanza municipal sobre reparto de agua existió y prevaleció. El estudio de
cuenca fue considerado necesario por el propio Estado en 1935.

\textbf{El problema de este territorio no es que falten herramientas. Es que las que hay se
usan de manera desigual, y que nadie está obligado a explicar por qué.}
\end{cautela}
